\section{Typing Cyclic Process Structures}
\label{sec:typing}

% PAR: section overview, typing reads the order, deadlock-free rules, coloured parts in sec:lock-freedom | example: ex-cycle | figure/rule: fig:typing-core | map: t-judgment

\subsection{Types and Context-Free Session Types}
\label{sec:types}

\begin{figure}[tp]
  \centering
  \maprow{\mapblock{ty-grammar}}
  \maprow{\mapblock{ty-session}}
  \maprow{\mapblock{ty-heads}}
  \caption{Types, context-free session types and active points}
  \label{fig:types}
  % FIG: map ty-grammar, ty-session, ty-heads (ty-dual, ty-payload in the appendix) | DF explained in sec:types | LF parts (credit, tickets, cost) explained in sec:lf-annotations
\end{figure}

% PAR: end type S_x, point in the subscript | example: ex-cycle | figure/rule: fig:types | map: ty-grammar
% PAR: arrow with upper point u and lower point l | example: ex-cycle | figure/rule: fig:types | map: ty-grammar, ty-heads
% PAR: boxes quantify over points and session variables under a constraint, existentials | example: ex-cs | figure/rule: fig:types | map: ty-grammar
% PAR: context-free session types, Skip, sequencing, mu, unfold, contractive, standard pointer Thiemann Vasconcelos | example: ex-ring | figure/rule: fig:types | map: ty-session, ty-equiv
% PAR: duality, flip of session variables | example: - | figure/rule: app:ty-dual | map: ty-dual
% PAR: payload binders, payload_at, pscoped | example: ex-ring | figure/rule: app:ty-payload | map: ty-payload
% PAR: active points ap, bounds bot and top | example: - | figure/rule: fig:types | map: ty-heads
% GAP: no map example uses sequencing S.S of context-free session types in a listing (ex-ring uses mu, ex-controller uses Skip in the lock-free view)
% GAP: contractivity example ex-sinst has no listing in the map, only in the Rocq

\subsection{The Typing Judgment}
\label{sec:judgment}

\begin{figure}[tp]
  \centering
  \maprow{\mapblock{t-judgment}}
  \maprow{\mapblock{t-contexts}}
  \maprow{\mapblock{t-side}}
  \caption{The typing judgment, contexts and side conditions on points}
  \label{fig:judgment}
  % FIG: map t-judgment, t-contexts, t-side | DF explained in sec:judgment | LF parts (bound N, effect, aging) explained in sec:lf-annotations
\end{figure}

% PAR: judgment Gamma, Delta, future F, world W | example: ex-cycle | figure/rule: fig:judgment | map: t-judgment
% PAR: future F, what the context does afterwards, F extended by the active points of the rest | example: ex-cycle | figure/rule: fig:judgment | map: t-judgment, t-contexts
% PAR: fact theory, hypotheses from boxes and existentials, Ord rules | example: ex-cs | figure/rule: Ord-World, Ord-Fact, Ord-Trans | map: t-contexts
% PAR: side conditions u < F, x < H, l below H | example: ex-cycle | figure/rule: fig:judgment | map: t-side
% PAR: without side conditions plain GV and the reversed cycle types | example: ex-cycle | figure/rule: fig:judgment | map: t-side

\subsection{Typing Rules}
\label{sec:rules}

\begin{figure}[tp]
  \centering
  \maprow{\maprule{E-Abs}}
  \maprow{\maprule{E-App}}
  \maprow{\maprule{E-Case}}
  \maprow{\maprule{E-Send}}
  \maprow{\maprule{E-Close}}
  \maprow{\maprule{E-Link}}
  \caption{Selected typing rules for functions, data and actions (all rules in \cref{app:typing})}
  \label{fig:typing-core}
  % FIG: map rules E-Abs, E-App (r-fun), E-Case (r-data), E-Send, E-Close, E-Link (r-act), E-Pair, E-Recv, E-Select, E-Branch, E-Spawn in the appendix | DF explained in sec:rules | LF parts (credits, effects, refill, link tickets) explained in sec:lf-annotations
\end{figure}

% PAR: E-Abs, body under the future of u, l below the captured ends | example: - | figure/rule: E-Abs | map: r-fun
% PAR: E-App, call waits at most at u, u < F | example: - | figure/rule: E-App | map: r-fun
% PAR: E-Pair, left component typed under the future extended by the right | example: - | figure/rule: app:r-data | map: r-data
% PAR: E-Send, E-Recv, E-Close, act at x, x < F, payload above x | example: ex-cycle | figure/rule: E-Send, E-Close | map: r-act, t-consts
% PAR: choice in the core, select to the dual field, branch to fields | example: ex-point-reuse | figure/rule: app:r-act | map: r-act, syn-moves
% PAR: E-Case types both arms in one world, dual points make both orders typable | example: ex-bit | figure/rule: E-Case | map: r-data
% PAR: E-Link, dual types at dual points | example: ex-copycat | figure/rule: E-Link | map: r-act
% PAR: E-Spawn, thunk | example: ex-cycle | figure/rule: app:r-threads | map: r-threads
% PAR: subtyping and structural rules, pointer to appendix | example: - | figure/rule: app:typing | map: t-subs

\begin{example}[The bit program \rocq{bit_deadlock_free}]
\label{ex:bit}
% Generated from map example ex-bit (The bit program). Do not edit.
\leavevmode\par\nobreak\noindent
\begin{minipage}{\linewidth}
\begin{maplisting}
main = \lstkw{let} ⟨e, e′⟩ = \lstkw{new} \lstkw{in}          \lstcmt{-- e at ρ, with the six facts below}
       \lstkw{let} ⟨c, c′⟩ = \lstkw{new} \lstkw{in}          \lstcmt{-- at the field ρ·0}
       \lstkw{let} ⟨d, d′⟩ = \lstkw{new} \lstkw{in}          \lstcmt{-- at the field ρ·1}
       \lstkw{let} ∗ = \lstkw{spawn} (λu. \lstkw{close} (\lstkw{send} ⟨inj₁ ∗, e′⟩)) \lstkw{in}
       \lstkw{let} ∗ = \lstkw{spawn} (λu. \lstkw{close} c′) \lstkw{in} \lstkw{let} ∗ = \lstkw{spawn} (λu. \lstkw{close} d′) \lstkw{in}
       \lstkw{let} ⟨v, e⟩ = \lstkw{recv} e \lstkw{in} \lstkw{let} ∗ = \lstkw{close} e \lstkw{in}
       \lstkw{case} v \lstkw{of} \lstkw{inl} w ⇒ \lstkw{let} ∗ = w \lstkw{in} \lstkw{let} ∗ = \lstkw{close} c \lstkw{in} \lstkw{close} d
               | \lstkw{inr} w ⇒ \lstkw{let} ∗ = w \lstkw{in} \lstkw{let} ∗ = \lstkw{close} d \lstkw{in} \lstkw{close} c
\end{maplisting}
\end{minipage}

% PAR: both arms in one world, six facts from the creation of e, one point per channel would give a cycle | example: ex-bit | figure/rule: E-Case, E-New | map: r-data, r-new
\end{example}

% Listing in the appendix (app:ex-copycat), cut from the body for space.
% PAR: five facts, forwarder at independent points, ring of two forwarders rejected, forwarder in a cut | example: ex-copycat | figure/rule: app:ex-copycat, E-Send, E-Recv, P-Connect | map: r-act, p-typing

\subsection{Creation and Recursion}
\label{sec:creation}

\begin{figure}[tp]
  \centering
  \maprow{\mapblock{r-new}}
  \maprow{\maprule{E-Box}}
  \maprow{\maprule{E-Force}}
  \maprow{\maprule{E-Open}}
  \caption{Typing rules for creation, boxes and existentials}
  \label{fig:typing-orders}
  % FIG: map r-new (E-New with acyc_new), E-Box, E-Force (r-boxes), E-Open (r-ex), E-Pack in the appendix | DF explained in sec:creation | LF parts (new_ok_N, tickets split m1 + m2 <= N, costs) explained in sec:lf-annotations and sec:lf-ranked
\end{figure}

% PAR: E-New, fresh base with a constraint or a field of a held point | example: ex-factory | figure/rule: E-New | map: r-new
% PAR: acyc_new, acyclic in the new tree, conservative over the context | example: ex-cycle | figure/rule: E-New | map: r-new, r-ex
% PAR: exits, facts from the new base down to the context | example: ex-gennew | figure/rule: E-New | map: r-new
% PAR: boxes, recursion, infinitely many facts in one constraint | example: ex-cs | figure/rule: E-Box, E-Force | map: r-boxes
% PAR: creation inside a box at a field of an argument point | example: ex-between | figure/rule: E-New, E-Box | map: r-new, r-boxes
% PAR: existentials abstract the points of a fresh base, E-Open conservative | example: ex-factory | figure/rule: E-Open, app:r-ex | map: r-ex
% PAR: one base and its fields, Milner scheduler, 18 facts | example: ex-milner | figure/rule: E-New | map: r-new
% PAR: recursive fork-join, consistent world that is not well-founded, typable here | example: ex-forkjoin | figure/rule: E-New, E-Box | map: r-new, w-wf

\begin{example}[Client and server cycle \rocq{cs_deadlock_free}]
\label{ex:cs}
% Generated from map example ex-cs (Client and server cycle). Do not edit.
\leavevmode\par\nobreak
\begin{maplisting}
main = let ⟨a, a′⟩ = new in                    -- a at ρ
       let ⟨b, b′⟩ = new in                    -- b at ρ·0, facts C at (ρ, ρ·0, ρ·1)
       let ∗ = spawn (λu. let ∗ = u in (force srv) ⟨a′, b′⟩) in
       (force cli) ⟨a, b⟩
cli  = !(f = λp. let ⟨a, b⟩ = p in let a = send ⟨∗, a⟩ in
                let ⟨u, b⟩ = recv b in let ∗ = u in (force f) ⟨a, b⟩)
srv  = !(g = λp. let ⟨a, b⟩ = p in let ⟨u, a⟩ = recv a in let ∗ = u in
                let b = send ⟨∗, b⟩ in (force g) ⟨a, b⟩)
-- both boxes : □∀ι κ ζ. C. ∗ ⊸ (S_ι ⊗ S′_κ ⊸^ζ ∗)
-- C = {ι+k ≺ (κ+k)‾, κ+k ≺ (ι+k+1)‾, ι+k ≺ ζ, κ+k ≺ ζ | k}
\end{maplisting}

% PAR: two loops over the same channels, infinite box constraint, next points at each call | example: ex-cs | figure/rule: E-Box, E-Force, E-New | map: r-boxes, r-new
\end{example}

% Listing in the appendix (app:ex-factory), cut from the body for space.
% PAR: one creation, two typings, existential at a fresh base or field of a held end | example: ex-factory | figure/rule: app:ex-factory, E-New, E-Pack | map: r-new, r-ex

\begin{example}[Milner's cyclic scheduler \rocq{ms_deadlock_free}]
\label{ex:milner}
% Generated from map example ex-milner (Milner's cyclic scheduler). Do not edit.
\leavevmode\par\nobreak\noindent
\begin{minipage}{\linewidth}
\begin{maplisting}
sendclose r = \lstkw{close} (\lstkw{send} ⟨∗, r⟩)
recvclose r = \lstkw{let} ⟨u, r⟩ = \lstkw{recv} r \lstkw{in} \lstkw{let} ∗ = u \lstkw{in} \lstkw{close} r
A₀ = \lstkw{let} ∗ = sendclose c₀ \lstkw{in} \lstkw{let} ∗ = recvclose d₀ \lstkw{in} sendclose a₀   \lstcmt{-- main}
Aᵢ = \lstkw{let} ∗ = recvclose dᵢ \lstkw{in} \lstkw{let} ∗ = sendclose aᵢ \lstkw{in} sendclose cᵢ   \lstcmt{-- i = 1, 2, spawned}
Pᵢ = recvclose bᵢ                                                 \lstcmt{-- i = 0, 1, 2, spawned}
\lstcmt{-- c₀, d₁ at ρ   a₁, b₁ at ρ·1   c₁, d₂ at ρ·2}
\lstcmt{-- a₂, b₂ at ρ·3   c₂, d₀ at ρ·4   a₀, b₀ at ρ·5}
\end{maplisting}
\end{minipage}

% PAR: six channels at one base and its fields, facts rise with the field label | example: ex-milner | figure/rule: E-New | map: r-new
\end{example}

\begin{example}[Recursive fork-join with replies \rocq{fj_deadlock_free}]
\label{ex:forkjoin}
% Generated from map example ex-forkjoin (Recursive fork-join with replies). Do not edit.
\leavevmode\par\nobreak\noindent
\begin{minipage}{\linewidth}
\begin{maplisting}
main = \lstkw{let} ⟨w, w′⟩ = \lstkw{new} \lstkw{in}                      \lstcmt{-- w at ρ, facts H at \ensuremath{\bar{ρ}}}
       \lstkw{let} ∗ = \lstkw{spawn} (λu. (\lstkw{force} wk) w′) \lstkw{in} \lstkw{close} w
wk   = !(f = λw. \lstkw{let} ⟨c₁, c₁′⟩ = \lstkw{new} \lstkw{in} \lstkw{let} ⟨c₂, c₂′⟩ = \lstkw{new} \lstkw{in}   \lstcmt{-- at ι·0 and ι·1}
             \lstkw{let} ∗ = \lstkw{spawn} (λu. (\lstkw{force} f) c₁′) \lstkw{in}
             \lstkw{let} ∗ = \lstkw{spawn} (λu. (\lstkw{force} f) c₂′) \lstkw{in}
             \lstkw{let} ∗ = \lstkw{close} c₁ \lstkw{in} \lstkw{let} ∗ = \lstkw{close} c₂ \lstkw{in} \lstkw{close} w)    \lstcmt{-- reply after both}
\end{maplisting}
\end{minipage}

% PAR: infinite descending chain, consistent, deadlock free, revisited in sec:lf-ranked | example: ex-forkjoin | figure/rule: E-New | map: r-new, w-wf
\end{example}

\subsection{Typing Processes}
\label{sec:proc-typing}

\begin{figure}[tp]
  \centering
  \maprow{\mapblock{p-typing}}
  \caption{Process typing}
  \label{fig:process-typing}
  % FIG: map p-typing (P-Term, P-Connect, P-Connect^n) | DF explained in sec:proc-typing | LF parts (N, eta, ranked) explained in sec:lf-annotations and sec:lf-ranked
\end{figure}

% PAR: P-Term, one thread, restricted channels, consistent world | example: ex-close | figure/rule: P-Term | map: p-typing
% PAR: P-Connect, cut on any set of channels, joined world consistent | example: ex-close | figure/rule: P-Connect, fig:worlds | map: p-typing, w-ext
% PAR: joined world checked as a whole, three components consistent in pairs | example: ex-ring3 | figure/rule: P-Connect | map: p-typing, w-ext
% PAR: private constants, renaming | example: ex-shared | figure/rule: P-Connect | map: p-typing, th-fund-proc
% PAR: static rejection of the flag, lemma on two closes | example: ex-flag | figure/rule: lem:two-closes | map: t-lemmas, p-typing

\begin{lemma}[Two closes, renaming \rocq{cl2_needs}]
\label{lem:two-closes}
\[ \mapblock{t-lemmas} \]
\end{lemma}

% Listing in the appendix (app:ex-ring3), cut from the body for space.
% PAR: consistent two at a time, inconsistent together | example: ex-ring3 | figure/rule: app:ex-ring3, P-Connect | map: p-typing, w-ext
